\section{Research assessment and the scalar distinction}
The immediate goal is positive integral occupancy of the original
Erd\H{o}s--Straus divisor atlas, not the construction of another object that
encodes occupancy after its answer is known. The preceding complexes built
from $D_\alpha=R/\gcd(R,g_\alpha)$ do retain failed states, their incoming
relations, and their supported zeros. Their formal contraction estimates,
however, also hold for arbitrary prescribed arrays of failed defects. Without
an additional arithmetic statement, those estimates do not force a successful
state. The observation-Gram and prime-indexed Koszul constructions also used
``support'' differently from the owner's original scalar distinction.

Here $A$ is a coefficient ring and
\[
 G(A)=\{\tau\}\sqcup A^\bullet,\qquad e=0_A^\bullet\ne\tau.
 \tag{S1}
\]
The original source distinguishes an empty face, a supported zero amplitude,
and an arbitrary supported amplitude. Its mixed-support monad permits a
previously absent coordinate to be present with residue zero; the source's
ordinary residue map still has $\{\tau\}$, not a supported infinitesimal fibre,
over $\tau$ \cite[\S6.1--6.2]{Owner}. Its actual internal quotient retains the
receiving support and the original boundary primitive \cite[\S1--4]{Derived}.
We use these distinctions, not a new three-element scalar system.

\section{Which inverses preserve absent and supported zero}
Let $\Gamma$ be a finite group and $A$ a nonzero commutative unital ring.
A coefficient vector in $G(A)[\Gamma]$ is exactly a pair
\[
 (S,f),\qquad S\subseteq\Gamma,\quad f\in A[\Gamma],\quad
 \operatorname{nz}(f)\subseteq S.
 \tag{S2}
\]
A coefficient outside $S$ is $\tau$, and a coefficient in $S$ whose amplitude
is zero is $e$. The complete inverse to (S2) reads those two cases separately.
Addition and convolution are
\[
 (S,f)+(T,g)=(S\cup T,f+g),\qquad
 (S,f)(T,g)=(ST,fg).
 \tag{S3}
\]
Indeed a product coefficient is present precisely when some present pair
contributes to it. Such contributions can cancel in amplitude but cannot
become absent. The global multiplicative identity is
$(\{1_\Gamma\},[1_\Gamma])$; the global additive identity is $(\varnothing,0)$.
Forgetting $S$ is an amplitude homomorphism with complete fibre
$\{S:S\supseteq\operatorname{nz}(f)\}$. It is not injective.

\begin{theorem}[Sparse split units and the nonnegative inverse obstruction]
The global units of $G(A)[\Gamma]$ are exactly the monomials
$a^\bullet[g]$ with $a\in A^\times$, $g\in\Gamma$.
Moreover, if an ordinary real group-algebra element and its inverse both have
nonnegative coefficients, both are monomials.
\end{theorem}
\begin{proof}
If $(S,f)(T,g)$ is the global identity, $ST=\{1_\Gamma\}$ and neither set is
empty. Fixing $t\in T$, injectivity of $s\mapsto st$ forces $|S|=1$; then
$|T|=1$. Their amplitudes multiply to one. This also proves the converse.
For ordinary nonnegative coefficients, the nonzero coefficient support of a
product is exactly the product of the two nonzero supports; the same argument
applies.
\end{proof}
This theorem concerns the declared global free coefficient semiring. It does
\emph{not} prohibit ordinary linear inverses in a fixed full-support fibre of
the SplitZero reconstruction. Those inverses have different domains and identities.

\paragraph{An exact six-cycle comparison.}
The ordinary target in the earlier source is $D_k=[3]+2[3-k]$ on $C_6$;
its Fourier transform is everywhere nonzero \cite[Abstract]{Center}.
At $k=2$, with additive indices modulo six, put
\[
 D=[3]+2[1],\qquad B=\frac{[3]-2[1]+4[5]}9.
\]
Ordinary multiplication gives $DB=[0]$. With its sparse supported lift,
\[
 \boxed{DB=1^\bullet[0]+e[2]+e[4]=E_H,\quad H=\{0,2,4\}.}
 \tag{S4}
\]
The six individual products give coefficients $1,0,0$ at $0,2,4$ respectively;
all three sites were used. Thus $E_H^2=E_H$, but $E_H$ is not the global identity.
Also
\[
 E_HD=[3]+2[1]+e[5]=D^{\rm sat},\quad
 E_HD^{\rm sat}=D^{\rm sat},\quad D^{\rm sat}B=E_H.
 \tag{S5}
\]
The target has an inverse relative to the displayed saturated identity after
retaining the new present-zero site. Its ordinary Fourier-unit assertion was
correct; it did not prove an inverse with the original sparse presence mask.
For any subgroup $H$, $(H,[1])$ gives this same idempotent completion:
it preserves amplitudes and replaces $S$ by $HS$. Keeping the old mask gives
the inverse as a decorated correspondence; forgetting it loses absence data.
Activation with amplitude zero is not an arithmetic witness.

\paragraph{Fourier observation.}
In an ordinary finite Fourier matrix and its inverse every entry is nonzero.
Under split matrix evaluation, any nonempty input presence mask therefore
becomes the full output mask, and inverse Fourier evaluation returns full
presence with the original amplitudes. This follows entry by entry from (S3)
and the ordinary Fourier inverse. In particular the composite is support
completion, not the identity on sparse split vectors. On a predeclared
full-support fibre the usual inverse is valid. No new Fourier theorem is
claimed; the change of carrier is the point.

\section{Original arithmetic before the answer is encoded}
Fix $R\equiv3\pmod4$, $a>R/2$, $\gcd(a,R)=1$, and put
\[
 n=4a-R>R,\quad a=\prod_{i=1}^kq_i^{e_i},\quad g_i=q_i\bmod R.
 \tag{A1}
\]
Primality of $n$ is not needed for the following constructions. The prime
applications separately check it. The primes $q_i$ are distinct and ordered,
and their multiplicities are never replaced by a generated subgroup.
Let $\Gamma_R=(\mathbb Z/R\mathbb Z)^\times$. Define
\[
 P^E_{g,e}=\sum_{f=0}^{2e}[g]^f,\qquad
 P^M_{g,e}=\sum_{\beta=-e}^{e}[g]^\beta,
 \quad t_E=-4^{-1},\quad t_M=-1\quad\text{in }\Gamma_R.
 \tag{A2}
\]
Then the coefficient of $t_c$ in $\prod_iP^c_{g_i,e_i}$ is exactly the
number of original channel-$c$ states. In E, $u=\prod q_i^{f_i}$; in M,
$u=a\prod q_i^{\beta_i}$. The target conditions are respectively
$R\mid4u+1$ and $R\mid u+a$. This explains why the two targets in (A2) are
constant when $R$ is fixed, without deleting the original $a$.

All coefficients in (A2) are nonnegative integer counts. Sending zero count
to $\tau$ and positive count $m$ to $m^\bullet$ is a semiring homomorphism
from $\mathbb N[\Gamma_R]$ to $G(\mathbb Q)[\Gamma_R]$, since no cancellation
or nontrivial zero divisor is possible. This is a declared sparse counting
carrier. Declaring every site present instead gives another valid carrier,
with the added zero sites retained. Neither is confused with evaluation of a
\emph{given} gate: its zero residue is the supported scalar $e$. Thus an absent
target, a signed cancellation, and a successful gate evaluation are different
observations with different fibres.

\paragraph{Exact positive denominators.}
For a target word, let $d=\gcd(a,u)$ and
\[
 h=d^2/u,\quad r=u/d,\quad s=a/d.
\]
At a prime $q$ with valuations $v_q(a)=e$, $v_q(u)=f$, the valuations of
$r,s,h$ are $(f-e)^+,(e-f)^+,e-|f-e|$. They are nonnegative, so
$a=hrs$, $u=hr^2$, and $\gcd(r,s)=1$. The ordered return is
\[
 E:\ \kappa=(nr+s)/R,\quad(a,hs\kappa,nhr\kappa),\qquad
 M:\ \lambda=(r+s)/R,\quad(a,nhs\lambda,nhr\lambda).
 \tag{A3}
\]
Reduction modulo $R$ and $4a\equiv n$ show that the displayed quotients are
positive integers. Multiplying the reciprocal sum by $nhrs\kappa$ gives
$n\kappa+nr+s=(n+R)\kappa=4hrs\kappa$ in E; the middle computation gives
$n\lambda+r+s=(n+R)\lambda=4hrs\lambda$. Hence both solve ES.
For ordered denominators $(a,Y,Z)$ the inverses are
\[
 u=\frac{a^2}{RY-na}\ (E),\qquad
 u=\frac{na^2}{RY-na}\ (M).
 \tag{A4}
\]
The channel, order, and complete valuation vector remain part of the return.
The Type I/II framework is inherited \cite[\S2, Propositions 2.2 and 2.6]{ET}.

\section{Exact infinite exponent arithmetic with positive coefficients}
The next construction starts from the prime factors in (A1), not from the
already computed defect or success projector.

\begin{theorem}[Positive rational carriers and full coefficient fibres]
In $\mathbb Q[\Gamma_R][[z]]$,
\[
 \boxed{\sum_{e\ge0}P^E_{g,e}z^e
  =\frac{1+z[g]}{(1-z)(1-z[g]^2)},\qquad
 \sum_{e\ge0}P^M_{g,e}z^e
  =\frac{1+z}{(1-z[g])(1-z[g]^{-1})}.}
 \tag{P1}
\]
Both right sides have specified positive geometric expansions. Each coefficient
has a bijection with the original bounded exponent choices.
\end{theorem}
\begin{proof}
For E expand the two geometric series and the two-term numerator. Their
indices $(j,k,\varepsilon)\in\mathbb N^2\times\{0,1\}$ give
$e=j+k+\varepsilon$, $f=2k+\varepsilon$. Conversely
$\varepsilon=f\bmod2$, $k=(f-\varepsilon)/2$,
$j=e-k-\varepsilon$. These are nonnegative exactly on $0\le f\le2e$.
For M the corresponding map is
$e=j+k+\varepsilon$, $\beta=j-k$. Its inverse is
$\varepsilon=(e-\beta)\bmod2$,
$j=(e-\varepsilon+\beta)/2$, $k=(e-\varepsilon-\beta)/2$;
nonnegativity is exactly $|\beta|\le e$.
\end{proof}
The rational equalities are in the ordinary formal group algebra. The
\emph{positive expansions} have canonical sparse split lifts. The denominator
expressions are not asserted to have global inverses in $G(\mathbb Q)[\Gamma_R][[z]]$.
This distinction is precisely (S4), now at infinite formal degree. Every
coefficient is a finite sum with its original labelled terms, so no divergent
summation or choice of analytic branch is involved.

\begin{theorem}[Periodic count formula and exact finite capacity reduction]
Let $o=\operatorname{ord}_{\Gamma_R}(g)$ and write $e=r+ok$, $0\le r<o$.
For either channel,
\[
 \boxed{P^c_{g,e}=P^c_{g,r}+2k\sum_{h\in\langle g\rangle}[h].}
 \tag{P2}
\]
In particular its coefficient support is unchanged by replacing $e$ with
$\min(e,\lfloor o/2\rfloor)$. That capacity is the least one at which the
whole subgroup is filled. Counts are not unchanged by this support reduction.
\end{theorem}
\begin{proof}
Increasing $e$ by $o$ adds $2o$ consecutive exponents in E. In M it adds
$o$ exponents at each end. Each added block of length $o$ traverses the
subgroup once, proving (P2) by iteration. A consecutive interval of length
$2e+1$ covers the cyclic subgroup exactly once its length reaches $o$.
Below $\lfloor o/2\rfloor$ that length is smaller than $o$. Nonnegative
convolution preserves equality of supports, so the reduction holds independently
in every factor of a product.
\end{proof}
Thus every fixed finite residue alphabet has an explicit multiaffine
quasipolynomial count on each residue class of its exponent parameters modulo
their orders: multiply the linear expressions (P2). This formula keeps arbitrarily
large multiplicities with finitely many group coordinates.

If several \emph{distinct} original primes have the same residue $g$, their
coefficient supports combine as one interval with capacity
$E_g=\sum_{q\equiv g}e_q$. Their counts instead retain
\[
 C_g(T)=\prod_{q\equiv g}(1+T+\cdots+T^{2e_q})
 \tag{P3}
\]
in E, and the analogous centered polynomial in M. The full inverse fibre of
aggregate exponent $f$ is all $(f_q)$ with $0\le f_q\le2e_q$ and
$\sum f_q=f$. Equivalently choose the coordinates successively and retain the
interval bounds imposed by the remaining capacities. These fibres are nonempty
for every $0\le f\le2E_g$, by induction on the number of intervals.
For two factors the fibre length is
$\min(2e_1,f)-\max(0,f-2e_2)+1$. One prime of multiplicity two has five choices;
two distinct primes of multiplicity one have nine. Their residue supports agree
but their original amplitudes do not.

\section{Complete local classification at residual 27}
Here $\Gamma_{27}=\langle2\rangle$ has order 18. The two target logarithms
are 7 in E and 9 in M. For $1\le j\le17$ put
\[
 E_j(a)=\sum_{q\equiv2^j\ (27)}v_q(a),\qquad
 c_j=\left\lfloor\frac{18/\gcd(j,18)}2\right\rfloor.
 \tag{C1}
\]
The identity-residue primes do not alter target existence; their coefficient
polynomials and factors in $a$ remain retained.

\begin{theorem}[Complete finite conductor at $R=27$]\label{thm:27}
There are exactly 174 coordinatewise minimal nonnegative capacity profiles
for which E or M has a target. They comprise 9 one-class profiles, 99 two-class
profiles and 66 three-class profiles. No minimal profile uses more classes.
Their maximum total capacity is nine.

For every $a>27/2$ coprime to 27, a channel hit exists if and only if its
profile $(E_j(a))$ dominates one of those 174 profiles. Consequently any hit
has a witness obtained by reserving at most nine \emph{actual prime occurrences}
in $a$, belonging to at most three nonidentity residue classes modulo 27.
All surplus factors remain in the original $a$ and ordered reconstruction.
\end{theorem}
\begin{proof}[Proof, including the finite certificate boundary]
Equation (P3) identifies coefficient \emph{support} with the aggregate interval
model. Equation (P2) reduces every capacity to $0\le E_j\le c_j$, preserving
target existence. Occupancy is increasing in every capacity. These reductions
apply to every exponent, not just a computed range.

The following complete finite calculation identifies its minimal occupied
profiles. Its full output and target words are in \texttt{catalogue27.json}.
First enumerate all $2^{17}=131072$ zero/one supports. A support $S$ has
18-bit masks $V_E(S),V_M(S)$, with the empty mask $\{0\}$. If $j\in S$ is
removed and $W$ is the preceding mask, its update is
\[
 V_E=W\cup(W+j)\cup(W+2j),\qquad
 V_M=W\cup(W+j)\cup(W-j),
 \tag{C2}
\]
with residues modulo18. Test membership of 7 and 9. This leaves 360 misses,
256 with only even logarithms and 104 with at least one odd logarithm.
A profile supported on even logarithms can never reach either odd target at
any multiplicity. For each of the other 104 supports enumerate its positive
capacities $1\le E_j\le c_j$, a total of 14040 profiles, using the same interval
updates. Keep an occupied profile precisely when all its immediate coordinate
predecessors miss. Together with the 85 squarefree minimal profiles this gives
174, with arity counts $9,99,66$ and maximum capacity sum nine. Both the masks
and the exact finite loop bounds are reproduced by the standalone verifier;
each retained profile includes a target word and every immediate predecessor
is separately checked.

An independent checker uses original unit residues, not logarithmic bit masks.
It starts at the zero profile and explores every unoccupied profile in the
capped grid. It visits 62812 such profiles and 629841 upward edges. Every
occupied boundary edge (567030 in total) dominates a listed profile whose
witness and minimality have already been checked in original residues.
Thus every monotone path either remains unoccupied or first crosses a checked
occupied boundary. This independently verifies the complete antichain, rather
than merely corroborating selected entries. It is a finite computer-assisted
proof of the displayed small classification; its general reduction was proved
before enumeration.

Finally, for a dominated minimal profile reserve its indicated number of
actual occurrences in each residue class. Allocate its target exponent to
those primes within their reserved capacities, using the interval fibres
(P3); this is always possible. E uses the resulting raw exponents. M uses
those centered exponents, with every unreserved prime at its original central
exponent. Equations (A3)--(A4) give the original marked integer witness.
\end{proof}

\paragraph{A short counterfactual consequence.}
The 104 squarefree misses containing an odd logarithm have arity distribution
$7,45,43,9$ at arities $1,2,3,4$. Hence, if the prime divisors of $a$ occupy
at least five distinct nonidentity residue classes modulo27 and at least one
is $2\bmod3$, the residual-27 shell necessarily has a hit. This criterion is
pointwise and uses only the given factorization, not an already-found divisor.
If such an odd-logarithm shell misses and has exactly four nonidentity classes,
its class set must be one of
\[
\begin{gathered}
 \{5,10,19,23\},\ \{5,8,11,13\},\ \{8,11,13,25\},\\
 \{2,4,7,17\},\ \{5,10,14,19\},\ \{10,14,19,23\},\\
 \{2,4,14,22\},\ \{2,4,14,17\},\ \{2,4,7,14\}.
\end{gathered}
\]
Further multiplicity restrictions are exactly the 174-profile avoidance test;
these nine support sets alone are necessary, not sufficient, for failure.
The alternative with every prime divisor $1\bmod3$ may contain arbitrarily
many occurrences and occupies only even logarithms, so it genuinely misses
this fixed residual. No global contradiction is claimed from this alternative.

The nine-occurrence bound is sharp for the local residue-profile model:
for the sole residue $23=2^{11}$, the E word requires raw exponent
$17\bmod18$, and M requires centered exponent $9\bmod18$.
Both first become available at capacity nine. No assertion that the
corresponding minimal prime instance has been found is needed.

Three residue classes really are necessary at a hard prime. At
\[
 p=19489\equiv169\pmod{840},\quad a=4879=7\cdot17\cdot41,\quad R=27,
\]
the complete 54-state shell has one hit:
\[
 E:\quad u=580601=7^2 17^2 41,\quad
 (h,r,s,\kappa)=(41,119,1,85896),
\]
\[
 (X,Y,Z)=(4879,3521736,8167578435576).
 \tag{C3}
\]
Every proper class subbudget misses both targets. The centred vector of this
E witness is $(1,1,0)$; the third class nevertheless occurs in its \emph{raw}
exterior target word. The arity statement refers to the declared E-raw/M-centered
budget model, not to the number of nonzero entries in every alternative chart.
Trial division through 139 proves primality. The catalogue, residues and
ordered inverse are all checked exactly.

\section{A hand-verifiable infinite-family subtheorem}
\begin{theorem}[All multiplicities for the $1,7,11$ residue family]
Suppose every prime factor of $a$ is $1,7$, or $11$ modulo 27. Let
$A=\sum_{q\equiv7}v_q(a)$ and $B=\sum_{q\equiv11}v_q(a)$. Then
\[
 \boxed{
 E\ne\varnothing\iff
 B\ge7\ \lor\ (A\ge1,\ B\ge3)\ \lor\ (A\ge2,\ B\ge1),}
 \tag{T1}
\]
\[
 \boxed{
 M\ne\varnothing\iff
 B\ge9\ \lor\ (A\ge1,\ B\ge5)\ \lor\ (A\ge2,\ B\ge1).}
 \tag{T2}
\]
Additional residue classes preserve sufficiency of these patterns but not
necessity. No assumption of primality for $n=4a-27$ is required.
\end{theorem}
\begin{proof}
The logarithms of 7 and 11 to base2 are $-2,-5$ modulo18. E is therefore
$2i+5j\equiv11\pmod{18}$ with $0\le i\le2A$, $0\le j\le2B$.
If $B=0$, parity excludes the target. If $A=0$, $j\equiv13\pmod{18}$,
so $B\ge7$. For $A=1$, the values $i=0,1,2$ require least nonnegative
$j=13,9,5$; hence $B\ge3$. If $A\ge2$, $(i,j)=(3,1)$ suffices as soon
as $B\ge1$. These exhaust all cases.

M is $2r+5s\equiv9\pmod{18}$ with $|r|\le A$, $|s|\le B$.
Again $B=0$ fails. At $A=0$ the least possible $|s|$ is9. At $A=1$
the three possibilities $r=-1,0,1$ have minimum possible $|s|$ respectively
$5,9,5$. At $A\ge2$, $(r,s)=(2,1)$ suffices when $B\ge1$.
The bounded composition fibres (P3) restore actual distinct prime factors;
residue1 factors do not change these congruences.
\end{proof}

The known empty shell $p=3361,a=847=7\cdot11^2$ has $(A,B)=(1,2)$.
Its failure is an inherited regression, not a new ES counterexample
\cite[\S2]{Boolean}. The theorem explains its exact threshold failure.
Adding one actual factor7 gives
\[
 a'=5929=7^2 11^2,\quad p'=23689\equiv169\pmod{840}.
\]
Its three hits include
\[
 M:\quad u=11,\quad(h,r,s,\lambda)=(11,1,539,20),
 \quad (X,Y,Z)=(5929,2809041620,5211580).
 \tag{T3}
\]
Primality is checked by all possible divisors through153.
This is an actual arithmetic extension, not an abstract activation of $e$.
For $a'=ta$, $n'=t(n+R)-R$ and $\gcd(t,R)=1$, the old E state embeds by
$u'=u$ and the old M state by $u'=tu$. Both retain their exact denominator
defect, since their gates are respectively unchanged and multiplied by $t$.
The inverse on the marked image is immediate from $t$. All three new hits
lie outside these old images. This does not prove a same-prime descent:
$n'$ and $n$ are explicitly different.

As a check of genuinely unbounded multiplicity, for the symbolic
$a=7^{10^{12}}11^{10^{12}+1}$, (P2) calculates without expanding $a$ or its
box the one-channel mass and counts
\[
\begin{array}{c|r}
\text{mass}&4000000000008000000000003\\
E&222222222222555555555555\\
M&222222222222777777777778.
\end{array}
\]
No primality claim for the corresponding enormous $n$ is made.

\section{An exact infinite Dirichlet carrier for the same arithmetic}
For fixed $R$, define $\mathsf P_c(a)=\prod_{q^e\parallel a}P^c_{q\bmod R,e}$
for $(a,R)=1$, and
\[
 \mathscr F_c(s)=\sum_{(a,R)=1}\mathsf P_c(a)a^{-s}.
\]
These are locally finite formal Dirichlet series first. For $\sigma=\Re s>1$
they converge absolutely in the group-algebra coefficient $\ell^1$ norm, since
\[
 \sum_{(a,R)=1}\tau(a^2)a^{-\sigma}
 =\prod_{q\nmid R}\frac{1+q^{-\sigma}}{(1-q^{-\sigma})^2}
 \le\frac{\zeta(\sigma)^3}{\zeta(2\sigma)}.
 \tag{D1}
\]
Uniqueness of the prime-factorization indices and the coefficient bijections
in (P1) justify the exact Euler products. In particular no success indicator
has been inserted before constructing the series.

For a Dirichlet character $\chi\bmod R$, including imprimitive ones, put
$\zeta_R(s)=\prod_{q\nmid R}(1-q^{-s})^{-1}$. Applying $\chi$ to the
finite group coordinate gives
\[
 \boxed{
 \chi(\mathscr F_E(s))=
 \frac{\zeta_R(s)L(s,\chi)L(s,\chi^2)}{L(2s,\chi^2)},\qquad
 \chi(\mathscr F_M(s))=
 \frac{\zeta_R(s)}{\zeta_R(2s)}L(s,\chi)L(s,\chi^{-1}).}
 \tag{D2}
\]
For E, rewrite $1+\chi(q)z=(1-\chi(q)^2z^2)/(1-\chi(q)z)$ in (P1).
For M use $1+z=(1-z^2)/(1-z)$. These prove (D2) prime by prime in the
absolutely convergent half-plane. The $L$-function conventions and their
Euler products are standard \cite[25.15.1--25.15.4]{DLMF}; no primitive
replacement or missing Euler factor at a prime dividing $R$ is permitted.
Finite character inversion gives the generating series of actual target
counts by $|\Gamma_R|^{-1}\sum_\chi\overline{\chi(t_c)}\chi(\mathscr F_c)$.
This amplitude observation retains the full original series as its source;
it is not a support-faithful inverse by (S4).

There is an explicit uniform-in-$R$ tail. Prime-exponent comparison gives
$\tau(a^2)\le d_3(a)$, since $2e+1\le(e+1)(e+2)/2$.
Counting ordered factorizations yields
$\sum_{a\le T}d_3(a)\le T(1+\log T)^2$.
Writing $a^{-\sigma}=\sigma\int_a^\infty t^{-\sigma-1}dt$, summing
nonnegative terms, and dropping the subtracted initial count gives, for $X\ge1$,
\[
 \boxed{
 \sum_{a>X}\tau(a^2)a^{-\sigma}
 \le \sigma X^{1-\sigma}
 \left[\frac{(1+\log X)^2}{\sigma-1}
 +\frac{2(1+\log X)}{(\sigma-1)^2}
 +\frac2{(\sigma-1)^3}\right].}
 \tag{D3}
\]
This proves the claimed convergence and explicit tail. The Euler carrier
uses all positive $a$ coprime to $R$, including its finite nonphysical prefix.
The actual first-half positive-shell carrier is its restriction to $a>R/2$:
in the ordinary amplitude series one subtracts the explicitly retained finite
Dirichlet polynomial $\sum_{a\le R/2,(a,R)=1}\mathsf P_c(a)a^{-s}$.
In the split carrier this is restriction of the index set with the removed
prefix recorded, not cancellation relabelled as absence. The result concerns
the full infinite arithmetic array at fixed modulus, not a nonzero coefficient
at every index. Nor is it the Gamma/period kernel of the other programme.
Imposing that $4a-R$ is prime changes the coefficient subsequence and does
not preserve the asserted Euler product without another argument.

\section{Verification, provenance and the remaining pointwise obligation}
The general inverse obstruction, positive coefficient maps, periodic formulas,
capacity reduction, two-class theorem and Dirichlet identities have the proofs
above. The 174-pattern theorem has the stated exhaustive finite certificate,
checked by two independently implemented algorithms; this is not independent
mathematical review.

The full residual-27 classification was also applied to all 4519 primes at
most two million in the six stipulated hard classes modulo840. It matched
actual counts in every instance: 1214 occupied residual-27 shells, 1880 E
states and 1800 M states. This is not new prime coverage. An independent
small-modulus check includes all first-half shells of 44 primes $1\bmod4$
through500, with 80032 canonical candidates and 1043 hits. Every output
witness is checked through (A4). No assertion about other residuals is
inferred from these tables.

For a hypothetical false prime, the profile of $a=(p+27)/4$ must avoid all
174 exact patterns. Analogous fixed-residual classifications can be generated
from (P2), but $R=27$ alone is not a universal selector. The next arithmetic
step must connect different shells of the \emph{same} $p$, retaining
\[
 a_{R+4t}=a_R+t,\qquad \gcd(a_R,a_R+t)=\gcd(a_R,t).
\]
Neither the validity of these identities nor an invertible ordinary Fourier
matrix proves that one of those correlated profiles is occupied. The work
here stops before that implication. It replaces further defect-based
repackaging by actual occurrence-level arithmetic and identifies exactly why
support activation and amplitude inversion cannot supply the implication on
their own. No historical-priority, Lean-build, global ES, or RH claim is made.
